\chapter{Fixings et flux}\label{ch:m2:fixings-and-flows}

Entre décembre 2007 et janvier 2013, des traders euro-dollar de quatre des
plus grandes banques mondiales, membres d'un salon de discussion qu'ils appelaient «~The Cartel~», utilisaient
un langage codé pour coordonner leurs transactions autour de deux moments de la journée~: 13~h~15
à Londres, lorsque le taux de référence de la Banque centrale européenne était
relevé, et 16~h, lorsque WM/Reuters fixait le taux de référence servant à valoriser les
ordres de nombreux gros clients. Ils
se communiquaient mutuellement les ordres de leurs clients, négociaient ensemble au fixing, et retenaient
leurs cours acheteurs ou vendeurs pour protéger leurs positions réciproques. En mai 2015, les quatre
banques ont accepté de plaider coupables aux États-Unis et de payer plus de 2.5
milliards de dollars d'amendes pénales, après que les régulateurs de Londres et de Washington avaient infligé
cinq banques plus de 3~milliards de dollars au total en novembre 2014. Ce chapitre
explique ce qu'est un \omterm{def:m2:fixings-and-flows:fix}{fixing de référence}, pourquoi autant de devises changent de mains à ce moment, les
flux de fin de mois et d'expiration des options qui s'y rencontrent, les manquements
constatés et les réformes qui ont suivi.

\section{Le fixing de référence}

\begin{definition}[Fixing de référence, fenêtre de fixing, ordre au fixing]\label{def:m2:fixings-and-flows:fix}
Un \emph{fixing de référence}\index{fixing de référence} est un taux de change publié à un
moment fixe de la journée par un administrateur, à partir de transactions et de cotations observées sur une courte
période autour de cet instant, la \emph{fenêtre de fixing}\index{fenêtre de fixing}. Un
\emph{ordre au fixing}\index{ordre au fixing} est un ordre passé par un client pour acheter ou vendre un
montant de devises au fixing, quel qu'en soit le résultat~; le dealer qui
l'accepte garantit le taux et supporte le coût de son obtention.
\end{definition}

Le plus utilisé est le taux de clôture WM/Reuters de 4pm à Londres
(\cref{fig:m2:fixings-and-flows:clock}), qui, selon les termes des
régulateurs américains, sert à valoriser les ordres de nombreux gros clients et à valoriser
les produits dérivés dans le monde entier. Un gérant dont le portefeuille est valorisé par rapport à un
indice de référence calculé à ce taux (livre~1, chapitre~3) a intérêt à négocier
exactement à ce moment-là~: une transaction réalisée au fixing n'ajoute aucune tracking error. Jusqu'en
février 2015, le fixing était calculé à partir de transactions sur une fenêtre d'une minute autour de
16:00~; à la suite des recommandations du Financial Stability Board en septembre
2014, WM l'a élargie le 15 février 2015 à cinq minutes, de 2.5 minutes avant
l'heure à 2.5 minutes après. La BCE publie ses propres taux de référence en euros
vers 16:00, heure d'Europe centrale, à partir d'une procédure menée vers 14:10, et précise que
les utiliser pour des transactions est «~fortement déconseillé~».

\begin{omfigure}
\centering
\begin{tikzpicture}[font=\small,x=0.62cm]
\draw[thick,-{Stealth}] (0,0) -- (19,0) node[right,font=\footnotesize] {heure de Londres};
\foreach \h in {0,2,...,18} {\draw (\h,0.1) -- (\h,-0.1) node[below,font=\scriptsize] {\h:00};}
\draw[omDef,thick] (6,0) -- (6,0.8);
\node[omDef,font=\footnotesize,anchor=south,align=center] at (6,0.8) {Cut d'expiration de Tokyo\\(15:00 Tokyo)};
\draw[omDef,thick] (15,0) -- (15,-0.75);
\node[omDef,font=\footnotesize,anchor=north,align=center] at (15,-0.75) {Cut d'expiration de New York\\(10:00 New York)};
\draw[omProp,thick] (13.17,0) -- (13.17,2.25);
\node[omProp,font=\footnotesize,anchor=south east,align=right] at (13.17,2.0) {Procédure des taux de référence de la BCE\\(14:10 CET)};
\draw[omThm,thick] (16,0) -- (16,2.25);
\node[omThm,font=\footnotesize,anchor=south west,align=left] at (16.15,2.25) {WM/Reuters\\fixing de 4pm};
\fill[omThm!25] (15.958,0) rectangle (16.042,0.35);
\end{tikzpicture}

\omcaption{Les moments de référence de la journée, à l'heure de Londres, en hiver. Le fixing WM/Reuters
est relevé sur cinq minutes autour de 16:00~; les taux de référence de la BCE proviennent
d'une procédure menée vers 14:10, heure d'Europe centrale~; les options de change expirent le plus souvent à
10:00 à New York ou 15:00 à Tokyo. Les horloges de l'heure d'été décalent certains d'entre eux d'une
heure. Schéma indicatif.}\label{fig:m2:fixings-and-flows:clock}
\end{omfigure}

La position du dealer est délicate. Un client lui demande d'acheter 1~milliard d'euros au
fixing. Le dealer ne connaît pas le fixing à l'avance~; il doit acheter lui-même les euros,
et ses propres achats font monter le cours. S'il les achète uniformément pendant la
fenêtre, il paie en moyenne environ le fixing et ne gagne que sa commission. S'il
en achète une partie avant la fenêtre, il paie des cours plus bas avant que ses propres
achats n'aient fait bouger le marché, puis ses achats restants et ceux de tous les autres
font monter le fixing, et il vend au client au fixing plus élevé.

\begin{proposition}[Ce que rapporte la pré-couverture d'un ordre au fixing]\label{prop:m2:fixings-and-flows:prehedge}
Supposons que le dealer achète $Q$ avec un impact permanent linéaire $\lambda$ par unité, une part
$p$ avant la fenêtre et le reste uniformément pendant celle-ci, et que le fixing soit la
moyenne des prix sur la fenêtre. Alors, à partir d'un prix de départ $P_0$ et en l'absence d'autre
variation de prix,
\[
\text{fixing} = P_0 + \frac{\lambda Q\,(1+p)}{2}, \qquad
\text{coût moyen} = P_0 + \frac{\lambda Q}{2}, \qquad
\text{gain} = \frac{\lambda Q^2\,p}{2}.
\]
Le gain est payé par le client, dont le fixing est plus élevé de $\lambda Q p/2$ que
sans pré-couverture.
\end{proposition}

\begin{proof}
Avant la fenêtre, le dealer achète $pQ$ à des prix croissants de $P_0$ à
$P_0 + \lambda pQ$, soit en moyenne $P_0 + \lambda pQ/2$. Dans la fenêtre, il achète $(1-p)Q$
à des prix croissants de $P_0 + \lambda pQ$ à $P_0 + \lambda Q$, dont la moyenne est le
fixing, $P_0 + \lambda Q(1+p)/2$. Le coût moyen vaut $p(P_0 + \lambda pQ/2) +
(1-p)\bigl(P_0 + \lambda Q(1+p)/2\bigr) = P_0 + \lambda Q/2$, quel que soit $p$. Le
gain est égal à $Q$ fois l'écart.
\end{proof}

\begin{omfigure}
\begin{tikzpicture}
\begin{axis}[width=11cm,height=5cm,
  xlabel={minutes par rapport à 16:00},ylabel={pips au-dessus du départ},
  tick label style={font=\small},label style={font=\small},
  xmin=-10,xmax=5,ymin=-0.2,ymax=3.6,grid=major,grid style={gray!25},
  legend columns=3,legend style={font=\footnotesize,at={(0.5,-0.3)},anchor=north,draw=none,fill=none,/tikz/every even column/.append style={column sep=4pt}},
  table/col sep=comma]
\fill[omThm!10] (axis cs:-2.5,-0.2) rectangle (axis cs:2.5,3.6);
\node[font=\scriptsize,omThm] at (axis cs:0,3.35) {fenêtre de fixing};
\addplot[omDef,very thick] table[x=t,y=p0]{figdata/markets-2/17-fixings-and-flows/path.csv};
\addplot[omProp,very thick,dashed] table[x=t,y=p50]{figdata/markets-2/17-fixings-and-flows/path.csv};
\addplot[omThm,very thick,dotted] table[x=t,y=p100]{figdata/markets-2/17-fixings-and-flows/path.csv};
\draw[omDef,thin] (axis cs:-2.5,1.5) -- (axis cs:2.5,1.5);
\draw[omThm,thin] (axis cs:-2.5,3.0) -- (axis cs:2.5,3.0);
\legend{tout dans la fenêtre (fixing $+1.5$),moitié pré-couverte (fixing $+2.25$),tout pré-couvert (fixing $+3$)}
\end{axis}
\end{tikzpicture}

\omcaption{Le prix de l'EURUSD pendant qu'un dealer achète 1~milliard d'euros pour un
\omterm{def:m2:fixings-and-flows:fix}{ordre au fixing} d'un client, avec un impact permanent de 0.3 pip pour 100~millions d'euros,
lorsqu'il achète tout dans la fenêtre de cinq minutes, la moitié avant, ou tout avant. Les
traits fins marquent le fixing dans le premier et le dernier cas. Le coût moyen du dealer
vaut 1.5 pip dans chaque cas~; le fixing, lui, ne l'est pas. Illustratif~; données~: le
tutoriel du chapitre.}\label{fig:m2:fixings-and-flows:path}
\end{omfigure}

\begin{example}[Un milliard d'euros au fixing]\label{ex:m2:fixings-and-flows:order}
Avec $\lambda = 0.3$ pip pour 100~millions d'euros, l'achat de 1~milliard d'euros fait monter le
cours de 3 pips. Acheté régulièrement dans la fenêtre, le fixing se situe 1.5 pip au-dessus du point de départ, et
le coût moyen du dealer est le même. À moitié acheté avant la fenêtre, le fixing monte de
2.25 pips, et le dealer gagne 0.75 pip sur 1~milliard d'euros, soit 75\,000~USD~; tout acheté avant,
150\,000~USD. Le client paie les mêmes montants, davantage que sans pré-couverture.
\end{example}

\section{Flux de fin de mois et rééquilibrages d'indices}

\begin{definition}[Rééquilibrage de fin de mois]\label{def:m2:fixings-and-flows:monthend}
Le \emph{rééquilibrage de fin de mois}\index{rééquilibrage de fin de mois} désigne les
opérations que les investisseurs effectuent à la fin de chaque mois, en partie au fixing de 16~h, pour
rétablir les couvertures de change et les pondérations de portefeuille que les mouvements de cours du mois ont
perturbées.
\end{definition}

Deux mécanismes rendent ces flux de fin de mois prévisibles. Un investisseur étranger qui couvre
la moitié de ses actions en dollars constate, après un mois où elles ont progressé de 4~\%, que sa
couverture en couvre moins de la moitié~: il vend davantage de dollars à terme, soit la moitié de 4~\% de ses
avoirs. Un fonds qui détient 60~\% d'actifs américains et 40~\% d'actifs européens
constate, après un mois où les premiers ont progressé de 5~\% et les seconds reculé de 1~\%,
que le fonds détient trop des premiers~; il vend des actifs américains, et donc des dollars, et achète
des actifs européens. Ces deux flux sont connus en signe, et à peu près en taille, plusieurs jours
à l'avance (\cref{fig:m2:fixings-and-flows:hedgeflow}).

\begin{omfigure}
\begin{tikzpicture}
\begin{axis}[width=11cm,height=4.8cm,
  xlabel={Rendement des actions américaines sur le mois (\%)},ylabel={ventes de dollars (Md USD)},
  tick label style={font=\small},label style={font=\small},
  xmin=-6,xmax=6,ymin=-95,ymax=95,grid=major,grid style={gray!25},
  legend columns=3,legend style={font=\footnotesize,at={(0.5,-0.3)},anchor=north,draw=none,fill=none,/tikz/every even column/.append style={column sep=8pt}},
  table/col sep=comma]
\addplot[omDef,thick,dashed] table[x=ret,y=h25]{figdata/markets-2/17-fixings-and-flows/hedgeflow.csv};
\addplot[omThm,very thick] table[x=ret,y=h50]{figdata/markets-2/17-fixings-and-flows/hedgeflow.csv};
\addplot[omProp,thick,dotted] table[x=ret,y=h75]{figdata/markets-2/17-fixings-and-flows/hedgeflow.csv};
\legend{couvert à 25~\%,couvert à 50~\%,couvert à 75~\%}
\end{axis}
\end{tikzpicture}

\omcaption{Ventes de dollars de fin de mois (négatif~: achats) nécessaires pour maintenir la
couverture de 2~billions de dollars d'actions américaines à son ratio, selon le rendement des actions sur le mois et
le ratio de couverture. Le flux est connu en signe dès que le rendement du mois l'est.
Illustratif~; données~: le tutoriel du chapitre.}\label{fig:m2:fixings-and-flows:hedgeflow}
\end{omfigure}

\begin{example}[Deux flux de fin de mois]\label{ex:m2:fixings-and-flows:flows}
Des investisseurs détenant 2~billions de dollars d'actions américaines, couvertes à 50~\%, vendent
$0.5 \times 4\% \times 2$ billions $=$ 40~milliards de dollars à terme après une hausse de 4~\%. Un
fonds de 100~milliards de dollars réparti 60/40, après que la part américaine a gagné 5~\% et la
part européenne a perdu 1~\%, doit vendre 1.44~milliard de dollars d'actifs américains et acheter
la même valeur d'actifs européens.
\end{example}

Tout le monde peut faire ces estimations, si bien que les flux connus à l'avance peuvent être
anticipés sur le marché~: un dealer ou un hedge fund qui attend des ventes de dollars au fixing
peut vendre le premier, et le cours peut bouger avant le fixing plutôt qu'à celui-ci. Ce que
l'investisseur paie pour la certitude d'un taux de référence est donc incertain, ce
qui justifie de répartir les grandes opérations de fin de mois sur une période plus longue ou
entre plusieurs algorithmes d'exécution.

\begin{dated}{2026-09}{La taille du stock}\label{dat:m2:fixings-and-flows:size}
Les détentions d'actions américaines de sociétés par le reste du monde étaient valorisées à
22.2~billions de dollars au deuxième trimestre 2026 (Federal Reserve Financial
Accounts). Aucun chiffre officiel ne donne la part de ce montant couverte contre le risque de change,
et les ratios de couverture varient selon l'investisseur et le pays.
\end{dated}

\section{Échéances d'options et cuts}

\begin{definition}[Cut d'expiration des options]\label{def:m2:fixings-and-flows:cut}
Un \emph{cut d'expiration des options}\index{cut d'expiration des options} est l'heure de la journée à laquelle
expire une option de change~: le plus souvent 10~h, heure de New York (le cut de New York) ou 15~h
heure de Tokyo (le cut de Tokyo).
\end{definition}

Les grandes options qui expirent à un cut, avec des prix d'exercice proches du spot,
créent leurs propres flux. Un dealer acheteur de l'option la couvre en vendant quand le cours monte et en achetant
quand il baisse, ce qui maintient le cours près du prix d'exercice à mesure que l'échéance approche~; un dealer vendeur
fait l'inverse et amplifie les mouvements. Au cut, le delta de l'option passe brusquement à zéro ou
à un, et sa couverture doit être défaite ou complétée. Les traders suivent les listes de
grandes échéances proches du spot en vue du cut du jour~; les options elles-mêmes font l'objet du
\cref{ch:m2:the-fx-options-market}.

\section{Le scandale des fixings et ce qui a changé}

\begin{definition}[Banging the close]\label{def:m2:fixings-and-flows:banging}
Le \emph{banging the close}\index{banging the close} consiste à concentrer des transactions
aux moments qui fixent un prix de référence ou de règlement-livraison, afin de le faire
bouger dans un sens favorable à une position.
\end{definition}

Les constats des régulateurs décrivent le mécanisme de la
\cref{prop:m2:fixings-and-flows:prehedge} devenu collusion. En novembre
2014, la UK Financial Conduct Authority a infligé à Citibank, HSBC, JPMorgan Chase, la
Royal Bank of Scotland et UBS une amende de 1.11~milliard de livres sterling pour n'avoir pas contrôlé
leurs activités de change au comptant sur les devises du G10 entre janvier 2008 et octobre 2013, période
pendant laquelle des traders ont partagé des informations confidentielles sur les ordres des clients et tenté
de manipuler les cours, y compris en concertation avec des traders d'autres établissements~; la US
Commodity Futures Trading Commission a infligé aux mêmes cinq banques plus de 1.4~milliard de dollars
pour tentative de manipulation des indices de référence de change, principalement le fixing de 16~h. En
mai 2015, Citicorp, JPMorgan, Barclays et la Royal Bank of Scotland ont accepté de
plaider coupable à une entente antitrust américaine visant à fixer les prix et à truquer les enchères en euros et
en dollars, avec des amendes de 395~millions à 925~millions de dollars chacune.

\begin{omfigure}
\begin{tikzpicture}
\begin{axis}[width=11cm,height=4.8cm,
  xlabel={la part de l'ordre achetée avant la fenêtre, $p$},ylabel={profit du dealer (milliers USD)},
  tick label style={font=\small},label style={font=\small},
  xmin=0,xmax=1,ymin=-80,ymax=380,grid=major,grid style={gray!25},
  legend columns=2,legend style={font=\footnotesize,at={(0.5,-0.3)},anchor=north,draw=none,fill=none,/tikz/every even column/.append style={column sep=8pt}},
  table/col sep=comma]
\addplot[name path=hi,draw=none,forget plot] table[x=p,y=hi]{figdata/markets-2/17-fixings-and-flows/pnl.csv};
\addplot[name path=lo,draw=none,forget plot] table[x=p,y=lo]{figdata/markets-2/17-fixings-and-flows/pnl.csv};
\addplot[omThm!20,forget plot] fill between[of=hi and lo];
\addplot[omThm,very thick,mark=*,mark size=1.2pt] table[x=p,y=pnl]{figdata/markets-2/17-fixings-and-flows/pnl.csv};
\legend{profit espéré (bande~: $\pm$ un écart-type)}
\end{axis}
\end{tikzpicture}

\omcaption{Profit du dealer sur un \omterm{def:m2:fixings-and-flows:fix}{ordre au fixing} de 1~milliard d'euros, selon la part
achetée avant la fenêtre, avec l'impact du chapitre et un mouvement aléatoire de deux
pips (écart-type) entre la pré-couverture et la fenêtre. Le profit espéré, prélevé
sur le client, croît avec la pré-couverture, et son risque aussi.
Illustratif~; données~: le tutoriel du chapitre.}\label{fig:m2:fixings-and-flows:pnl}
\end{omfigure}

Le fixing a évolué en conséquence. La fenêtre élargie rend la moyenne plus difficile à déplacer,
et WM a ajouté des sources de cours. Le \omterm{def:m2:fx-spot-microstructure:code}{FX Global Code}
(\cref{ch:m2:fx-spot-microstructure}) indique qu'un dealer traitant des \omterm{def:m2:fixings-and-flows:fix}{ordres au fixing}
ne doit pas partager inconsidérément des informations ni chercher à influencer le taux,
qu'il s'agisse de collusion ou d'autres moyens, ne doit pas chercher à influencer intentionnellement le fixing
pour en tirer profit, et doit tarifer ces ordres de façon transparente et cohérente
avec le risque qu'il supporte. Son principe~11 n'autorise la pré-couverture que pour un dealer
agissant en principal, pour gérer le risque d'un ordre client anticipé d'une manière
conçue pour profiter au client, et exige que cette pratique soit expliquée aux
clients. Opérer en amont d'un \omterm{def:m2:fixings-and-flows:fix}{ordre au fixing} n'est donc ni interdit ni libre~:
la question décisive est de savoir s'il sert le client, et la frontière reste débattue.

\section{Tutoriel~: flux de fin de mois et ordre au fixing}\label{tut:m2:fixings-and-flows:flows}

\begin{tutorial}
\textbf{Objectif.} Estimer les flux de couverture et de rééquilibrage de portefeuille de fin de mois, et
chiffrer un \omterm{def:m2:fixings-and-flows:fix}{ordre au fixing}, avec et sans pré-couverture. \textbf{État final~:}
\cref{fig:m2:fixings-and-flows:path,fig:m2:fixings-and-flows:pnl},
\cref{ex:m2:fixings-and-flows:order,ex:m2:fixings-and-flows:flows} et les
valeurs du problème du week-end.
\begin{tutsteps}
\item \textbf{Les flux}~: couverture et rééquilibrage des pondérations.
\omcode{code/firm/fixflow/firm_fixflow.py}{14}{23}{Flux de couverture et rééquilibrage de portefeuille en fin de mois.}\label{lst:m2:fixings-and-flows:flows}
\item \textbf{L'\omterm{def:m2:fixings-and-flows:fix}{ordre au fixing}}~: la forme fermée donnée par
la \cref{prop:m2:fixings-and-flows:prehedge}, et une simulation avec un mouvement aléatoire.
\omcode{code/firm/fixflow/firm_fixflow.py}{26}{51}{Fixing, coût pour le dealer et profit d'un ordre au fixing.}\label{lst:m2:fixings-and-flows:fix}
\item \textbf{Exécuter} \texttt{fixflow\_demo.month\_end()},
\texttt{fixflow\_demo.fixing\_order} pour $p = 0$, 0.5 et 1, et
\texttt{fig\_fixflow.py}.
\end{tutsteps}
\textbf{Pour aller plus loin.} Ajouter un impact temporaire, qui s'estompe après chaque
transaction, et montrer qu'il rend l'achat dans la fenêtre plus coûteux que le fixing~;
puis trouver la part de pré-couverture qui minimise le risque du dealer
pour un profit espéré nul.
\end{tutorial}

\section{Construction~: l'estimateur des flux de fin de mois}\label{bld:m2:fixings-and-flows:fixflow}

\begin{build}
\bfield{Rôle} L'entreprise miniature traite les \omterm{def:m2:fixings-and-flows:fix}{ordres au fixing} des clients et les transactions
autour des fins de mois~: elle doit estimer les flux auxquels elle fera face et connaître, et
communiquer l'économie de la façon dont elle exécute les \omterm{def:m2:fixings-and-flows:fix}{ordres au fixing}. \bfield{Interface} \texttt{hedge\_rebalance(value\_usd, asset\_return,
hedge\_ratio)}~; \texttt{weight\_rebalance(values, targets)}~;
\texttt{fix\_and\_cost(q, lam, p, p0)}~; \texttt{dealer\_pnl(q, lam, p)}~;
\texttt{simulate\_pnl(q, lam, p, sigma\_pre, n, seed)}.
\bfield{Règles} Ventes en dollars positives~; impact permanent linéaire~; le fixing comme
moyenne des cours de la fenêtre~; le mouvement aléatoire appliqué après la pré-couverture.
\bfield{Tests d'acceptation} \texttt{code/firm/fixflow/tests/}~: signes du flux
de couverture~; le fonds vend sa zone surperformante~; sans pré-couverture, le dealer achète
au fixing~; le coût indépendant de $p$~; la pré-couverture ajoute du risque.
\bfield{Pour aller plus loin} Flux par pays, à partir des données publiques de détention et des
rendements mensuels des indices~; impact temporaire~; le fixing calculé à partir d'un carnet d'ordres simulé, comme le fait WM
en prenant la médiane des taux observés.
\end{build}

\begin{omsources}
\item US Department of Justice, ``Five major banks agree to parent-level guilty
pleas'', 20 mai 2015.
\item Financial Conduct Authority, communiqué de presse, 12 novembre 2014~; Commodity
Futures Trading Commission, communiqué de presse 7056-14, 12 novembre 2014.
\item Financial Stability Board, \emph{Foreign Exchange Benchmarks: Final
Report}, septembre 2014~; WM/Reuters presentation to the ECB FX Contact Group,
novembre 2015.
\item European Central Bank, euro foreign exchange reference rates~; Federal Reserve, Financial Accounts of the United States.
\end{omsources}

\section{Exercices}

\begin{exercise}[$\star$]\label{exo:m2:fixings-and-flows:1}
Expliquer pourquoi un fonds indiciel voudrait échanger des devises au fixing de 16~h
plutôt qu'en cours de journée, à un meilleur prix.
\end{exercise}

\begin{exercise}[$\star$]\label{exo:m2:fixings-and-flows:2}
Un investisseur japonais détient 50~milliards de dollars d'actions américaines couvertes à 70~\%. Elles
progressent de 3~\% sur le mois. Calculer le montant négocié en fin de mois.
\end{exercise}

\begin{exercise}[$\star$]\label{exo:m2:fixings-and-flows:3}
Expliquer ce qui a changé dans le fixing WM/Reuters en février 2015, et pourquoi il rend
la manipulation plus difficile.
\end{exercise}

\begin{exercise}[$\star\star$]\label{exo:m2:fixings-and-flows:4}
Dans l'\cref{ex:m2:fixings-and-flows:order}, donner le fixing et le profit du dealer s'il
pré-couvre 30~\%.
\end{exercise}

\begin{exercise}[$\star\star$]\label{exo:m2:fixings-and-flows:5}
Un fonds alloué à 60/40 entre les États-Unis et l'Europe, avec 100~milliards de dollars, voit sa part américaine baisser de 3~\%
et sa part européenne monter de 2~\% en dollars. Quelles opérations doit-il réaliser~?
\end{exercise}

\begin{exercise}[$\star\star$]\label{exo:m2:fixings-and-flows:6}
Un dealer est long d'une grosse option EURUSD dont le prix d'exercice est juste au-dessus du spot et qui expire
au cut de New York. Comment sa couverture affecte-t-elle le cours le matin~?
\end{exercise}

\begin{exercise}[$\star\star\star$]\label{exo:m2:fixings-and-flows:7}
\emph{Programmation.} Avec \texttt{simulate\_pnl}, donner la moyenne et l'écart-type
du profit du dealer pour $p = 0.5$ et $p = 1$, avec les paramètres du chapitre.
\end{exercise}

\begin{exercise}[$\star\star\star$]\label{exo:m2:fixings-and-flows:8}
\emph{Trouver l'erreur.} «~Le dealer achète avant la fenêtre, ce n'est que de la couverture~:
son coût moyen est le même quoi qu'il fasse, donc le client n'est pas lésé.~»
Corriger ce raisonnement.
\end{exercise}

\section{Problème~: Garanti au fixing}

\begin{problem}[{Problème du week-end --- ce que rapporte un ordre au fixing au dealer}]\label{pb:m2:fixings-and-flows:1}
Un fonds de retraite demande à un dealer d'acheter 1~milliard d'euros au fixing de 16~h en fin de mois.
Le dealer estime un impact permanent de 0.3~pip pour 100~millions d'euros, et un
mouvement aléatoire de deux pips (écart-type) entre 15~h~50 et la fenêtre.

\medskip\noindent\textbf{Partie I --- Sans pré-couverture.}
\begin{enumerate}
\item De combien l'achat du dealer fait-il bouger l'EURUSD~?
\item Donner le fixing, par rapport au point de départ.
\item Donner le coût moyen et le profit du dealer.
\item Quel risque le dealer supporte-t-il~?
\item Que facture le dealer pour le service, et comment~?
\end{enumerate}

\medskip\noindent\textbf{Partie II --- Pré-couverture.}
\begin{enumerate}[resume]
\item Le dealer achète la moitié avant 15~h~57~min~30~s. Donner le fixing.
\item Donner son profit espéré, et qui le paie.
\item Donner l'écart-type de son profit.
\item Donner la même chose pour une pré-couverture totale.
\item Pourquoi le coût moyen ne dépend-il pas de la part pré-couverte~?
\end{enumerate}

\medskip\noindent\textbf{Partie III --- Autres flux.}
\begin{enumerate}[resume]
\item D'autres clients vendent 600~millions d'euros au même fixing. Qu'est-ce qui change~?
\item Le rééquilibrage de couverture de fin de mois ajoute 40~milliards de dollars de ventes de dollars. Dans quel
sens pousse-t-il l'EURUSD, et quand~?
\item Pourquoi le cours pourrait-il bouger avant la fenêtre plutôt que pendant~?
\item Comment le fonds pourrait-il réduire son coût~?
\item Qu'exige le \omterm{def:m2:fx-spot-microstructure:code}{FX Global Code} du dealer~?
\end{enumerate}

\medskip\noindent\textbf{Partie IV --- Jugement.}
\begin{enumerate}[resume]
\item Quand la pré-couverture relève-t-elle d'une gestion légitime du risque~?
\item Qu'est-ce qui a rendu criminelle la conduite de 2008--2013~?
\item Pourquoi une fenêtre de cinq minutes aide-t-elle, et qu'est-ce qu'elle n'empêche pas~?
\item Énoncer le \emph{résultat nommé}~: le profit espéré du dealer sur l'ordre, avec
et sans pré-couverture.
\item En une phrase~: qui paie quand un dealer trade avant un \omterm{def:m2:fixings-and-flows:fix}{ordre au fixing}~?
\end{enumerate}
\end{problem}

\section{Questions d'entretien}

\begin{interviewq}[$\star$ \iqroles{trader, banque}]\label{iq:m2:fixings-and-flows:1}
Qu'est-ce que le fixing WM/Reuters, et qui l'utilise~?
\end{interviewq}

\begin{interviewq}[$\star$ \iqroles{trader, chercheur}]\label{iq:m2:fixings-and-flows:2}
Pourquoi les flux de change de fin de mois sont-ils prévisibles, et dans quel sens vont-ils après
un mois fort pour les actions américaines~?
\end{interviewq}

\begin{interviewq}[$\star\star$ \iqroles{trader}]\label{iq:m2:fixings-and-flows:3}
Un client vous confie un \omterm{def:m2:fixings-and-flows:fix}{ordre au fixing} de 2~milliards d'euros. Comment l'exécutez-vous, et que
dites-vous au client~?
\end{interviewq}

\begin{interviewq}[$\star\star$ \iqroles{chercheur}]\label{iq:m2:fixings-and-flows:4}
Montrer qu'une pré-couverture d'un \omterm{def:m2:fixings-and-flows:fix}{ordre au fixing} rapporte de l'argent au dealer dans un modèle à
impact linéaire. D'où vient cet argent~?
\end{interviewq}

\begin{interviewq}[$\star\star$ \iqroles{banque, trader}]\label{iq:m2:fixings-and-flows:5}
De quoi portaient les affaires de benchmark des changes de 2014 et 2015, et qu'est-ce qui a changé
en conséquence~?
\end{interviewq}

\begin{interviewq}[$\star\star\star$ \iqroles{développeur, chercheur}]\label{iq:m2:fixings-and-flows:6}
Concevoir un système de surveillance qui détecte le \omterm{def:m2:fixings-and-flows:banging}{banging the close} autour du fixing de 16~h.
\end{interviewq}
